\clearpage
\SourceLeaf{445}
\fancyhead[L]{\small STÖRUNGEN DER PALLAS. NACHLASS.}
\pdfbookmark[1]{Nachlass. [II.] Darstellung einer neuen Methode}{pallas-II-exposition}
\begin{center}
{\Large NACHLASS.\par}
\bigskip
{\large [II.] DARSTELLUNG\par}
{\large EINER NEUEN METHODE\par}
ZUR BERECHNUNG DER PLANETENSTÖRUNGEN\par
MIT ANWENDUNG AUF DIE NUMERISCHE BERECHNUNG\par
DER STÖRUNGEN DER BEWEGUNG DER PALLAS.
\end{center}
\begin{flushright}
\begin{minipage}{.94\textwidth}
\small\raggedright
. . . . . their motions, periods and their laws\\
Give me to scan.\quad \textsc{Thomsons} Autumn 1346.
\par\medskip
See, how associate round their central sun,\\
Their faithful rings the circling planets run\\
Exactly tracing their appointed sphere.\\
\hspace*{1em}Boyle on deity: Elegant Extracts p. 126.
\end{minipage}
\end{flushright}
\noindent\hspace*{\parindent}Die Berechnung der Planetenstörungen lässt sich in zwei verschiedenen Formen durchführen. In der einen Form bezieht man die Bewegung eines Planeten auf konstante oder nur säkularen Änderungen unterworfene Elemente. Den unter dieser Annahme berechneten Radiusvektor, die Länge in der Bahn und die Breite korrigiert man durch periodische Gleichungen. In der anderen Form bezieht man die Bewegung auf eine veränderliche Ellipse, die die Bewegung zu jedem Zeitpunkt ohne weitere Korrektur genau darstellt. Es ist hinlänglich bekannt, mit welchem Erfolg man sich bisher für die Störungen der früher bekannten Planeten ausschließlich der ersten Form bedient hat. Man weiß jedoch, dass dieser Erfolg vor allem von mehreren günstigen Umständen abhängt, die bei einigen der neuen Planeten nicht mehr in gleicher Weise vorliegen. Die Bahn der Pallas ist gegen die Bahn des Jupiter um 34 Grad geneigt und hat eine \SourceLeaf{446}Exzentrizität, die fast gleich \(\frac{1}{4}\) ist. Die Bahn der Juno ist bei sehr beträchtlicher Neigung noch exzentrischer. In diesen Fällen erschienen die älteren Methoden undurchführbar.
\par Die neuen Methoden, die den Gegenstand dieser Abhandlung bilden, liefern die Störungen unmittelbar in der zweiten Form. Es wäre in der Tat leicht, die einmal gefundenen Störungen der Elemente in Störungen der drei Koordinaten umzurechnen, wenn man dies für zweckmäßig hielte. Wir sind jedoch der Ansicht, dass man sich diese Arbeit ersparen und bei den Störungen der Elemente bleiben kann. Zwar gibt es sechs solcher Elemente, und die Störungen der drei Koordinaten scheinen bei numerischen Anwendungen nur die Hälfte der Arbeit zu erfordern. Dieser Nachteil entfällt jedoch immer dann, wenn mehrere Orte des Planeten zu berechnen sind, die durch einen mäßigen Zeitraum getrennt sind. In diesem in der Praxis häufigsten Fall genügt ein einziger Satz von Elementen, während die andere Methode die Berechnung der Ungleichheiten der Koordinaten für mindestens zwei verschiedene Epochen erfordert. Bezieht man alle Ungleichheiten auf die Elemente selbst, so muss man auf den Vorteil einer Hilfstafel für die Mittelpunktsgleichung und den Radiusvektor in einer unveränderlichen Ellipse verzichten. Dieser Vorteil fällt bei einem Planeten, der so vielen Ungleichheiten unterliegt wie Pallas, ebenfalls kaum ins Gewicht und kann sogar durch andere Vereinfachungen der Rechnung ausgeglichen werden, die hier anzuführen überflüssig wäre.
\par\bigskip\centerline{\rule{3cm}{.4pt}}\par
